\chapter{Power Sources and Batteries}
\label{ch:power-sources}

\section{Ideal and Real Sources}

An \textbf{ideal voltage source} maintains a fixed voltage regardless of
load current; an \textbf{ideal current source} maintains a fixed current
regardless of load voltage. Real sources (batteries, power supplies,
generators) approximate the ideal voltage source but have a small
\textbf{internal resistance} $R_{\text{int}}$ in series, which causes the
terminal voltage to sag under load:
\[
\boxed{U_{\text{terminal}} = \text{EMF} - I R_{\text{int}}}
\]
This single equation explains why a weak or discharged battery's voltage
recovers somewhat when a load is removed (the $IR_{\text{int}}$ drop
disappears) and why a heavily loaded supply's voltage is lower than its
open-circuit (no-load) rating.

\begin{examplebox}[Voltage sag under load]
A 12\,V battery with internal resistance 0.2\,$\Omega$ supplies 8\,A to a
transceiver. Find the terminal voltage.

\emph{Solution.}
\[
U_{\text{terminal}} = 12 - (8 \times 0.2) = 12 - 1.6 = 10.4\,\text{V}.
\]
This sag is one reason mobile/portable radio equipment often specifies a
minimum operating voltage below the nominal 12\,V or 13.8\,V rating.
\end{examplebox}

\section{Maximum Power Transfer}

For a source with fixed internal resistance $R_{\text{int}}$ driving a
variable load $R_L$, the power delivered \emph{to the load} is maximized
when $R_L = R_{\text{int}}$ (the \textbf{maximum power transfer
theorem}) --- at that point exactly half the source's total generated
power is delivered to the load and half is dissipated internally, so
this condition maximizes load power but not overall efficiency. This
principle underlies impedance-matching practice throughout RF systems
(Chapter~\ref{ch:transmission-lines}): matching source and load
impedance maximizes power transfer, though in power-supply and battery
contexts, efficiency (minimizing $R_{\text{int}}$ relative to the load)
is usually the more relevant goal instead.

\section{Cells and Batteries}

A single electrochemical \textbf{cell} converts stored chemical energy
into electrical energy; a \textbf{battery} is one or more cells combined
(in series for higher voltage, in parallel for higher capacity, or both).

\begin{center}
\renewcommand{\arraystretch}{1.2}
\begin{tabularx}{\linewidth}{@{}l l l l Y@{}}
\toprule
Chemistry & Cell V & Rechg.? & Energy density & Notes \\
\midrule
Alkaline (Zn-MnO$_2$) & 1.5\,V & No & Low & Cheap, common, poor at high current \\
NiMH & 1.2\,V & Yes & Moderate & Low self-discharge (LSD) types common \\
Lead-acid & 2.0\,V/cell & Yes & Low--mod. & High current capable; heavy \\
Li-ion & 3.6--3.7\,V & Yes & High & Requires protection circuitry \\
LiFePO$_4$ & 3.2\,V & Yes & High & Safer chemistry, flatter discharge curve \\
\bottomrule
\end{tabularx}
\end{center}

\subsection{Series and Parallel Battery Combinations}

\begin{keyidea}
\textbf{Cells in series} add voltage; capacity (in amp-hours) stays the
same as a single cell. \textbf{Cells in parallel} add capacity; voltage
stays the same as a single cell. A common 12\,V lead-acid battery is
internally six 2\,V cells in series.
\end{keyidea}

\begin{safetynote}[Mixing cells]
Never combine cells of different chemistry, capacity, or state of charge
in the same series or parallel string. Mismatched cells can be forced
into reverse charge, over-discharge, or unequal current sharing, any of
which can cause overheating, venting, or fire --- especially with
lithium-based chemistries, which additionally require dedicated
protection/balancing circuitry in multi-cell packs.
\end{safetynote}

\section{Amp-Hour Capacity and Runtime}

Battery capacity is rated in \textbf{amp-hours (Ah)} or milliamp-hours
(mAh) --- a direct application of $Q = It$ from
Chapter~\ref{ch:electricity}, expressed with $t$ in hours instead of
seconds. Approximate runtime for a given load current is
\[
t \approx \frac{\text{Capacity (Ah)}}{I_{\text{load}} (\text{A})}
\]
(approximate because usable capacity typically falls at high discharge
current, a behavior described empirically by \textbf{Peukert's law} for
lead-acid batteries, and because most chemistries should not be
discharged to 0\% without damage or reduced cycle life).

\begin{examplebox}[Estimating runtime]
A handheld transceiver draws 300\,mA average on receive and is powered
by a 2000\,mAh battery. Estimate the receive-only runtime.

\emph{Solution.}
\[
t \approx \frac{2000\,\text{mAh}}{300\,\text{mA}} \approx 6.7\,\text{hours}.
\]
Real-world runtime will be shorter due to periods of transmit current
(much higher than receive current) and capacity derating at higher
discharge rates.
\end{examplebox}

Energy capacity in watt-hours (useful for comparing different battery
voltages on an equal footing) is
\[
W_h = \text{Ah} \times U_{\text{nominal}}.
\]

\section{Charging Considerations}

Different chemistries require different charge algorithms (constant
current then constant voltage for Li-ion; different profiles for NiMH
and lead-acid), and using the wrong charger or charge rate can
permanently damage capacity, or in the case of lithium chemistries,
create a serious fire hazard. \textbf{C-rate} notation describes charge
or discharge current relative to capacity: a ``1C'' rate for a 2\,Ah
battery is 2\,A; a ``0.5C'' rate is 1\,A.

\section{DC Power Supplies}

A mains-powered DC \textbf{power supply} typically consists of a
transformer (Chapter~\ref{ch:transformers}) to change AC mains voltage
to a convenient level, a rectifier (diodes, Chapter~\ref{ch:semiconductors})
to convert AC to pulsating DC, a filter capacitor to smooth the
pulsations into a nearly steady DC voltage, and a voltage regulator to
hold the output constant despite input ripple and load changes. A
\textbf{linear regulator} dissipates the voltage difference between
input and output as heat (simple, low-noise, but inefficient at large
voltage drops); a \textbf{switching regulator} rapidly switches a
storage inductor or capacitor to convert voltage with much higher
efficiency, at the cost of potential RF noise/interference that must be
filtered --- a genuine concern for a radio amateur choosing a supply to
sit next to a sensitive receiver.

\begin{quickreview}
\item Real sources sag under load: $U_{\text{terminal}} = \text{EMF} -
IR_{\text{int}}$.
\item Maximum power transfer occurs when load resistance equals source
internal resistance.
\item Series cells add voltage; parallel cells add capacity; never mix
chemistries, capacities, or states of charge in a pack.
\item Capacity (Ah) combined with load current gives approximate runtime,
$t \approx \text{Ah}/I$; watt-hours allow comparison across voltages.
\item Linear regulators are simple and quiet but lossy; switching
regulators are efficient but can generate RF interference.
\end{quickreview}

\begin{practiceproblems}
\item A power supply has an open-circuit voltage of 13.8\,V and internal
resistance 0.05\,$\Omega$. Find the terminal voltage at a 20\,A load.
\item Four 1.2\,V, 2000\,mAh NiMH cells are connected in series. Find the
resulting pack voltage and capacity.
\item Estimate the runtime of a 7\,Ah battery supplying a constant 1.4\,A
load.
\item Explain, using $Q=It$, why amp-hours (not just amps) are the
correct way to compare battery ``size.''
\item A radio amateur reports RF hash/noise across several bands after
switching from a linear to a switching power supply. Explain a likely
cause and one practical mitigation.
\end{practiceproblems}
